\documentclass[11pt]{article}

\usepackage[a4paper,margin=0.72in]{geometry}
\usepackage{amsmath,amssymb}
\usepackage{booktabs}
\usepackage{tabularx}
\usepackage{array}
\usepackage{enumitem}
\usepackage{xcolor}
\usepackage{microtype}
\usepackage{fancyhdr}
\usepackage{listings}

\pagestyle{fancy}
\fancyhf{}
\lhead{DA3410}
\rhead{Viva Question Set-3}
\cfoot{\thepage}

\setlength{\parindent}{0pt}
\setlength{\parskip}{4pt}
\setlist[itemize]{leftmargin=1.4em,itemsep=1pt,topsep=2pt}
\setlist[enumerate]{leftmargin=1.6em,itemsep=1pt,topsep=2pt}

\lstset{
  language=Python,
  basicstyle=\ttfamily\small,
  columns=fullflexible,
  keepspaces=true,
  showstringspaces=false,
  frame=single,
  framerule=0.3pt,
  aboveskip=4pt,
  belowskip=4pt,
  xleftmargin=0.5em,
  xrightmargin=0.5em
}

% Student copy: answers and marking rubrics are intentionally omitted
% from the source, not merely hidden by a toggle.
\newcounter{question}
\newenvironment{question}{%
  \refstepcounter{question}
  \par\medskip
  \noindent\textbf{Q\thequestion.}\par\smallskip
}{\par\medskip}

\newcommand{}{\textit{Student answers orally; no code writing required.}}

\begin{document}

\begin{center}
    {\LARGE \textbf{DA3410 Viva Questions Set-3}}\\[4pt]
\end{center}

% Bank Q4
\begin{question}
A student wants to fit a scaler or PCA before evaluating on validation data.
Which data should be used to fit the scaler/PCA?
\end{question}

% Bank Q8
\begin{question}
Give the activation output and local derivative for the input below.
Use the convention that the ReLU derivative at zero is $0$.

\begin{center}
\begin{tabular}{lc}
\toprule
Activation and input & Required output \\
\midrule
Leaky ReLU, $x=-2$, $\alpha=0.1$ & $f(x)$ and $f'(x)$ \\
\bottomrule
\end{tabular}
\end{center}
\end{question}

% Bank Q11
\begin{question}
There is one bug in this implementation. What is it? 

\begin{lstlisting}
def sigmoid(x):
    return 1 / 1 + torch.exp(-x)
\end{lstlisting}
\end{question}

% Bank Q15
\begin{question}
Identify the problem in the code and explain what is wrong with the resulting
validation measurement.

\begin{verbatim}
for xb, yb in val_loader:
    pred = forward(xb, params)
    loss = loss_fn(pred, yb)
    backward(loss)
    update(params, lr)
\end{verbatim}
\end{question}

% Bank Q18
\begin{question}
A manual autodiff system performs the forward operation shown below.
For the local backward function, assume the upstream gradient and only the
quantities explicitly listed in \texttt{saved} are available.

Is the saved information sufficient to compute the local backward pass?
If not, state what additional forward quantity must be retained.

\begin{verbatim}
y = a * b
saved = {"a": a}
\end{verbatim}
\end{question}

% Bank Q20
\begin{question}
For the values below, give $\partial L/\partial a$ and
$\partial L/\partial b$. 

\begin{lstlisting}
y = a * b
# upstream gradient is g = dL/dy
\end{lstlisting}

\begin{center}
\begin{tabular}{ccc}
\toprule
$a$ & $b$ & $g$ \\
\midrule
3 & 2 & 1 \\
\bottomrule
\end{tabular}
\end{center}
\end{question}

% Bank Q24
\begin{question}
For \texttt{z = x @ W}, which candidate correctly computes the gradient with
respect to \texttt{x}? 

\begin{lstlisting}
# Candidate A
grad_x = grad_z @ W.T

# Candidate B
grad_x = W.T @ grad_z

# Candidate C
grad_x = grad_z.T @ W
\end{lstlisting}
\end{question}

% Bank Q26
\begin{question}
How many batches are produced and what is the size of the last batch?


\begin{lstlisting}
for i in range(0, N, batch_size):
    xb = X[i:i + batch_size]
\end{lstlisting}

\begin{center}
\begin{tabular}{cc}
\toprule
$N$ & \texttt{batch\_size} \\
\midrule
8 & 4 \\
\bottomrule
\end{tabular}
\end{center}
\end{question}

% Bank Q29
\begin{question}
The validation loss improves, but the supposedly saved ``best'' parameters
keep changing during later training. Identify the bug. 

\begin{lstlisting}
if val_loss < best_loss:
    best_loss = val_loss
    best_params = params
\end{lstlisting}
\end{question}

% Bank Q31
\begin{question}
The following code is used to initialize a deep network before measuring
layerwise gradient magnitudes.

State whether the setup is intended to produce \textbf{stable},
\textbf{vanishing}, or \textbf{exploding} gradient behaviour.

\begin{verbatim}
for l in range(depth):
    W[l] = he_init(fan_in, fan_out)
activation = relu
\end{verbatim}
\end{question}


\end{document}
